\subsection{取向 torsor 的超对称结构、wreath 跳跃律与寄存器下界}\label{subsec:conclusion-orientation-torsor-supermonoidal-wreath-gauge-register}
沿定理 \ref{thm:xi-binary-jump-orientation-unification} 的统一口径，\(\mathrm{Or}\) 同时编码离散 uplift 与解析分支跳跃。
本小节将该口径在张量、乘积、分层与寄存器复杂度方向上的终态结构律统一为可直接调用的命题簇。

\begin{proposition}[\texorpdfstring{\(\ZZ_2\)}{Z2}-torsor 的平方范式化与奇偶幂坍缩]\label{prop:conclusion-z2torsor-square-parity-collapse}
\leanverified{paper_conclusion_z2torsor_square_parity_collapse}
设 \(P\) 为 \(\ZZ_2\)-torsor，并以
\[
P\otimes Q:=(P\times Q)/\Delta(\ZZ_2)
\]
定义 Picard 张量结构。
则存在不依赖额外选择的自然同构
\[
\Theta_P:\ P\otimes P\ \xrightarrow{\ \sim\ }\ \ZZ_2.
\]
从而对任意整数 \(n\ge 1\) 有规范同构
\[
P^{\otimes n}\cong
\begin{cases}
\ZZ_2,& n\ \text{为偶数},\\
P,& n\ \text{为奇数}.
\end{cases}
\]
\end{proposition}
\begin{proof}
任取 \((p,q)\in P\times P\)。
由 torsor 作用自由且传递，存在唯一 \(\varepsilon\in\ZZ_2\) 使 \(q=\varepsilon\cdot p\)。
定义
\[
\delta(p,q):=\varepsilon.
\]
对任意 \(\eta\in\ZZ_2\)，有 \(\delta(\eta\cdot p,\eta\cdot q)=\delta(p,q)\)，故 \(\delta\) 下推为
\[
\bar\delta:(P\times P)/\Delta(\ZZ_2)\to\ZZ_2.
\]
\(\bar\delta\) 显然双射，且与 \(\ZZ_2\)-作用相容，得 \(\Theta_P\)。
奇偶幂结论由 \(\Theta_P\) 迭代直接推出。证毕。
\end{proof}

\begin{theorem}[交换群目标上的对称幺半跳跃理论仅由 \texorpdfstring{$2$}{2}-扭决定]\label{thm:conclusion-monoidal-jump-theory-a2torsion-classification}
\leanverified{paper\_conclusion\_monoidal\_jump\_theory\_a2torsion\_classification}
设 \(A\) 为交换群，并记
\[
A[2]:=\{a\in A:\ 2a=0\}.
\]
考虑对称幺半函子
\[
J:(\mathbf{FinSet}^{\cong},\sqcup)\to (\mathbf{Tors}(A),\otimes),
\qquad
J(\varnothing)\ \text{平凡}.
\]
则在对称幺半自然同构意义下：
\begin{enumerate}
  \item \(J\) 的等价类由唯一元素 \(a_J\in A[2]\) 完全刻画；
  \item \(a_J=0\) 当且仅当 \(J\) 与平凡函子等价；
  \item 对任意 \(a_J\in A[2]\)，记 \(\phi_{a_J}:\ZZ_2\to A\) 为唯一同态 \(1\mapsto a_J\)，并定义结构群延拓
  \[
  (\phi_{a_J,\ast}\mathrm{Or})(S)
  :=
  \mathrm{Or}(S)\times_{\ZZ_2,\phi_{a_J}} A.
  \]
  则 \(J\simeq \phi_{a_J,\ast}\mathrm{Or}\)。
\end{enumerate}
\end{theorem}
\begin{proof}
对每个 \(n\ge 0\) 取骨架对象 \([n]=\{1,\dots,n\}\)。
由于 \(J([n])\) 为 \(A\)-torsor，其自同构群自然同构于 \(A\)，故 \(J\) 在自同构群上诱导群同态
\[
\chi_n:S_n\to \Aut(J([n]))\cong A.
\]
对称幺半自然性给出块和可加律
\[
\chi_{m+n}(\sigma\oplus\tau)=\chi_m(\sigma)+\chi_n(\tau)
\qquad
(\sigma\in S_m,\ \tau\in S_n).
\]
令
\[
a:=\chi_2((12))\in A.
\]
由
\[
\chi_n\bigl((12)\oplus\mathrm{id}_{n-2}\bigr)=\chi_2((12))=a
\]
可知所有换位在 \(\chi_n\) 下的像都等于 \(a\)：因为任意换位都与 \((12)\oplus\mathrm{id}_{n-2}\) 共轭，而目标 \(A\) 交换。
再由换位平方为恒等，得
\[
2a=0,
\]
故 \(a\in A[2]\)。

又因 \(A\) 交换，\(\chi_n\) 必经由交换化
\[
S_n^{\mathrm{ab}}\cong \ZZ_2
\qquad (n\ge 2)
\]
因子分解。
因此存在唯一同态 \(\phi_a:\ZZ_2\to A\) 使
\[
\chi_n=\phi_a\circ \operatorname{sgn}
\qquad (n\ge 2).
\]
另一方面，由定理 \ref{thm:xi-time-part9i-monoidal-rigidity-by-two-points} 与命题 \ref{prop:bdry-orientation-functor-symmetric-monoidal-koszul}，
\(\mathrm{Or}\) 是 \((\mathbf{FinSet}^{\cong},\sqcup)\to(\mathbf{Tors}(\ZZ_2),\otimes)\) 的唯一非平凡二值跳跃原子，并携带对称幺半结构。
将其经 \(\phi_a\) 作结构群延拓后，得到的 \(A\)-torsor 值函子 \(\phi_{a,\ast}\mathrm{Or}\) 在每个 \([n]\) 上诱导的自同构表示恰为 \(\phi_a\circ\operatorname{sgn}\)，与 \(J\) 的 \(\chi_n\) 完全一致。

最后，\(\mathbf{FinSet}^{\cong}\) 是由单点对象生成的自由对称幺半群胚，而 \(\mathbf{Tors}(A)\) 为 Picard 群胚；
在单位对象已固定为平凡 torsor 的前提下，一旦各阶层的对称群作用 \(\chi_n\) 固定，对称幺半函子的等价类便唯一确定。
故 \(J\simeq \phi_{a,\ast}\mathrm{Or}\)。
取 \(a_J:=a\) 即得唯一分类。
当 \(a_J=0\) 时，\(\phi_{a_J}\) 为零同态，\(\phi_{a_J,\ast}\mathrm{Or}\) 即为平凡函子，故第 (2) 条成立。证毕。
\end{proof}

\begin{corollary}[奇素数阶跳跃与连续相位跳跃的禁制律]\label{cor:conclusion-no-oddprime-or-continuous-phase-monoidal-jump}
沿用定理 \ref{thm:conclusion-monoidal-jump-theory-a2torsion-classification} 的设定。
\begin{enumerate}
  \item 若 \(A[2]=0\)（例如 \(A=\ZZ\) 或 \(A=\ZZ/p\ZZ\) 且 \(p\) 为奇素数），则任意
  \[
  J:(\mathbf{FinSet}^{\cong},\sqcup)\to (\mathbf{Tors}(A),\otimes),
  \qquad
  J(\varnothing)\ \text{平凡},
  \]
  都与平凡函子等价。
  \item 若 \(A=U(1)\)，则任意此类 \(J\) 都经由其 \(2\)-扭子群
  \[
  U(1)[2]=\{\pm 1\}
  \]
  因子化；等价地，全对称自然性下不存在真正连续的相位跳跃，所有非平凡跳跃都塌缩为二值符号。
\end{enumerate}
\end{corollary}
\begin{proof}
第 (1) 条由定理 \ref{thm:conclusion-monoidal-jump-theory-a2torsion-classification} 中 \(a_J\in A[2]\) 的唯一分类立即得到。
第 (2) 条中，\(U(1)[2]=\{\pm1\}\)，故同一定理表明任意 \(J\) 必与某个
\(
\phi_{a_J,\ast}\mathrm{Or}
\)
等价，其中 \(\phi_{a_J}\) 的像包含于 \(\{\pm1\}\)。
于是所谓“连续相位”在该函子学口径下被强制坍缩为二值符号。证毕。
\end{proof}

\leanverified{paper\_conclusion\_orientation\_supermonoidal\_koszul}
\begin{theorem}[取向函子在不交并下的超对称单子结构]\label{thm:conclusion-orientation-supermonoidal-koszul}
对有限集 \(S,T\)，存在自然同构
\[
\mu_{S,T}:\ \mathrm{Or}(S)\otimes \mathrm{Or}(T)\xrightarrow{\ \sim\ }\mathrm{Or}(S\sqcup T),
\]
满足结合与幺元约束。
并且对交换同构 \(\tau_{S,T}:S\sqcup T\to T\sqcup S\) 有
\[
\tau_{S,T}^{\ast}\circ \mu_{S,T}
=
(-1)^{|S||T|}\,
\mu_{T,S}\circ \mathrm{swap}.
\]
等价地，\(\mathrm{Or}\) 在 \(\ZZ_2\)-分次意义下为超对称 monoidal，其 braiding 精确由 Koszul 因子 \((-1)^{|S||T|}\) 给出。
\end{theorem}
\begin{proof}
由命题 \ref{prop:bdry-orientation-functor-symmetric-monoidal-koszul} 可得对应结论（该命题采用 \(\mathrm{Or}(S\sqcup T)\to \mathrm{Or}(S)\otimes\mathrm{Or}(T)\) 的等价取向）。
将其取逆即得本式方向。
块交换奇偶性为 \((-1)^{|S||T|}\)，故得到 Koszul 因子。证毕。
\end{proof}

\begin{theorem}[笛卡尔积下的 Kronecker--取向分解]\label{thm:conclusion-orientation-cartesian-kronecker}
设有限集 \(S,T\) 的基数分别为 \(d,e\)。
则存在自然同构
\[
\Phi_{S,T}:\ \mathrm{Or}(S\times T)\xrightarrow{\ \sim\ }\mathrm{Or}(S)^{\otimes e}\otimes \mathrm{Or}(T)^{\otimes d}.
\]
\end{theorem}
\begin{proof}
命题 \ref{prop:bdry-orientation-cartesian-product-exponent-law} 已给出同一公式。
其证明本质上由行列式恒等式
\(
\det(U\otimes W)=\det(U)^{\dim W}\det(W)^{\dim U}
\)
与取向行列式实现（命题 \ref{prop:bdry-orientation-torsor-determinant-realization-jumpclass}）组成。
这正对应于标准 wreath 嵌入下的 Kronecker 分解机制。证毕。
\end{proof}
\leanverified{paper\_conclusion\_orientation\_cartesian\_kronecker}

\begin{corollary}[积取向的二次坍缩与偶度屏蔽]\label{cor:conclusion-orientation-cartesian-parity-collapse}
在定理 \ref{thm:conclusion-orientation-cartesian-kronecker} 的同构下，结合命题 \ref{prop:conclusion-z2torsor-square-parity-collapse} 得
\[
\mathrm{Or}(S\times T)\ \cong\
\mathrm{Or}(S)^{\otimes (e\bmod 2)}
\otimes
\mathrm{Or}(T)^{\otimes (d\bmod 2)}.
\]
因此：
\begin{enumerate}
  \item \(d,e\) 皆偶时，\(\mathrm{Or}(S\times T)\cong \ZZ_2\)；
  \item \(d\) 偶、\(e\) 奇时，\(\mathrm{Or}(S\times T)\cong \mathrm{Or}(S)\)；
  \item \(d\) 奇、\(e\) 偶时，\(\mathrm{Or}(S\times T)\cong \mathrm{Or}(T)\)；
  \item \(d,e\) 皆奇时，\(\mathrm{Or}(S\times T)\cong \mathrm{Or}(S)\otimes\mathrm{Or}(T)\)。
\end{enumerate}
\end{corollary}
\leanverified{paper\_conclusion\_orientation\_cartesian\_parity\_collapse}
\begin{proof}
将定理 \ref{thm:conclusion-orientation-cartesian-kronecker} 中幂次按奇偶化简即可。证毕。
\end{proof}

\begin{theorem}[有限笛卡尔积的取向 torsor 仅由奇偶型决定]\label{thm:conclusion-orientation-finite-cartesian-parity-type}
\leanverified{paper\_conclusion\_orientation\_finite\_cartesian\_parity\_type}
设 \(r\ge 2\)，\(S_1,\dots,S_r\) 为有限集合，并记
\[
d_i:=|S_i|
\qquad
(1\le i\le r).
\]
则存在自然同构
\[
\mathrm{Or}\!\Bigl(\prod_{i=1}^r S_i\Bigr)
\cong
\bigotimes_{i=1}^r
\mathrm{Or}(S_i)^{\otimes \prod_{j\ne i}d_j}.
\]
特别地，\(\mathrm{Or}(\prod_{i=1}^r S_i)\) 仅依赖于 \((d_1,\dots,d_r)\) 的奇偶型，并满足精确分类
\[
\mathrm{Or}\!\Bigl(\prod_{i=1}^r S_i\Bigr)
\cong
\begin{cases}
\ZZ_2,& \text{若至少有两项 }d_i\text{ 为偶数},\\[2mm]
\mathrm{Or}(S_{i_0}),& \text{若恰有一项 }d_{i_0}\text{ 为偶数},\\[2mm]
\bigotimes_{i=1}^r\mathrm{Or}(S_i),& \text{若全部 }d_i\text{ 为奇数}.
\end{cases}
\]
\end{theorem}
\begin{proof}
对 \(r=2\) 时，第一式即为定理 \ref{thm:conclusion-orientation-cartesian-kronecker}。
现对 \(r\) 作归纳，并设结论已对 \(r-1\) 个集合成立。
记
\[
S':=\prod_{i=1}^{r-1}S_i.
\]
由定理 \ref{thm:conclusion-orientation-cartesian-kronecker}，
\[
\mathrm{Or}(S'\times S_r)
\cong
\mathrm{Or}(S')^{\otimes d_r}\otimes \mathrm{Or}(S_r)^{\otimes |S'|}.
\]
由归纳假设与 Picard 张量结构的结合律，
\[
\mathrm{Or}(S')^{\otimes d_r}
\cong
\bigotimes_{i=1}^{r-1}
\mathrm{Or}(S_i)^{\otimes d_r\prod_{\substack{1\le j\le r-1\\ j\ne i}}d_j}
=
\bigotimes_{i=1}^{r-1}
\mathrm{Or}(S_i)^{\otimes \prod_{j\ne i}d_j},
\]
而
\[
\mathrm{Or}(S_r)^{\otimes |S'|}
=
\mathrm{Or}(S_r)^{\otimes \prod_{j<r}d_j}
=
\mathrm{Or}(S_r)^{\otimes \prod_{j\ne r}d_j}.
\]
合并即得第一式。
\par
再由命题 \ref{prop:conclusion-z2torsor-square-parity-collapse}，\(P^{\otimes n}\) 只依赖于 \(n\bmod 2\)。
故仅当 \(\prod_{j\ne i}d_j\) 为奇数时，第 \(i\) 个因子才保留。
若至少有两项 \(d_i\) 为偶数，则对每个 \(i\) 都有 \(\prod_{j\ne i}d_j\) 为偶数，因而整体平凡；
若恰有一项 \(d_{i_0}\) 为偶数，则仅 \(\prod_{j\ne i_0}d_j\) 为奇数；
若全部 \(d_i\) 为奇数，则所有因子均保留。
证毕。
\end{proof}

\begin{theorem}[wreath 结构下的符号加法与偶度屏蔽]\label{thm:conclusion-wreath-sign-additivity-even-shielding}
令
\[
W:=\mathfrak S_d^{\,e}\rtimes \mathfrak S_e
\hookrightarrow
\mathfrak S_{de}
\]
为标准 wreath 嵌入，并把符号同态写作加法型
\(
\mathrm{sgn}_n:\mathfrak S_n\to\ZZ_2
\)。
则对任意 \(\big((\sigma_1,\dots,\sigma_e),\tau\big)\in W\) 有
\[
\mathrm{sgn}_{de}\!\big((\sigma_1,\dots,\sigma_e),\tau\big)
=
\Big(\sum_{j=1}^{e}\mathrm{sgn}_d(\sigma_j)\Big)+d\cdot \mathrm{sgn}_e(\tau)
\quad\text{于 }\ZZ_2.
\]
特别地：
\begin{enumerate}
  \item \(d\) 偶时，外层置换 \(\tau\) 对总二值跳跃不可见；
  \item \(d\) 奇时，总跳跃等于“块内跳跃之和 + 外层跳跃”。
\end{enumerate}
\end{theorem}
\begin{proof}
对应置换矩阵可分解为
\[
P=(P_\tau\otimes I_d)\cdot \mathrm{diag}(P_{\sigma_1},\dots,P_{\sigma_e}),
\]
故
\[
\det(P)=\det(P_\tau)^d\cdot\prod_{j=1}^{e}\det(P_{\sigma_j}).
\]
将行列式取 \(\ZZ_2\) 符号即得公式。两条推论是对 \(d\) 奇偶的直接代入。证毕。
\end{proof}

\begin{proposition}[fold-gauge 的中心二扭与二值电荷格分离]\label{prop:conclusion-foldbin-gauge-center-charge-gap}
设
\[
G_{\mathrm{bin},m}:=\mathcal{G}^{\mathrm{bin}}_m
=\prod_{x\in X_m}\mathfrak S_{d_m(x)},
\]
并定义
\[
s_m:=\#\{x\in X_m:\ d_m(x)\ge 2\},\qquad
t_m:=\#\{x\in X_m:\ d_m(x)=2\}.
\]
则有：
\begin{enumerate}
  \item \(\mathrm{Hom}(G_{\mathrm{bin},m},\ZZ_2)\cong (\ZZ_2)^{s_m}\)；
  \item \(Z(G_{\mathrm{bin},m})[2]\cong (\ZZ_2)^{t_m}\)；
  \item 存在自然坐标嵌入
  \[
  Z(G_{\mathrm{bin},m})[2]\hookrightarrow \mathrm{Hom}(G_{\mathrm{bin},m},\ZZ_2),
  \]
  其像恰为由二点纤维坐标生成的子空间；因此“非中心二值方向”的维数为 \(s_m-t_m\)。
\end{enumerate}
\end{proposition}
\begin{proof}
前两条分别由定理 \ref{thm:window6-foldbin-gauge-abelianization-even-parity} 与命题 \ref{prop:window6-foldbin-gauge-center-vs-charge-separation} 的结构分解直接给出。
嵌入由“二点纤维中心换位 \(\mapsto\) 该纤维符号字符”逐坐标定义；其单射性与像描述同样见命题 \ref{prop:window6-foldbin-gauge-center-vs-charge-separation} 的坐标化证明。
维数差即 \(s_m-t_m\)。证毕。
\end{proof}

\begin{theorem}[等变 \texorpdfstring{\(a\)}{a}-分裂的最小寄存器态下界]\label{thm:conclusion-equivariant-a-splitting-register-lower-bound}
\leanverified{paper\_conclusion\_equivariant\_a\_splitting\_register\_lower\_bound}
令 \(\Delta\) 为 \(k\)-元集合，\(1\le a\le k-1\)，群 \(\mathfrak S_k\) 以自然方式作用于 \(\Delta\)。
称 \((R,A)\) 为一个 \(\mathfrak S_k\)-等变 \(a\)-分裂机制，若：
\begin{enumerate}
  \item \(R\) 是有限 \(\mathfrak S_k\)-集；
  \item \(A\subseteq \Delta\times R\) 在对角作用下不变；
  \item 对每个 \(r\in R\)，纤维 \(A_r:=\{x\in\Delta:(x,r)\in A\}\) 满足 \(|A_r|=a\)。
\end{enumerate}
则：
\begin{enumerate}
  \item 上述数据等价于一个 \(\mathfrak S_k\)-等变映射
  \[
  \varphi:R\to \binom{\Delta}{a},
  \qquad
  A=\{(x,r):x\in\varphi(r)\};
  \]
  \item 必有
  \[
  |R|\ge \binom{k}{a};
  \]
  \item 下界可达：取 \(R=\binom{\Delta}{a}\) 与 \(\varphi=\mathrm{id}\) 即得。
\end{enumerate}
特别地，当 \(a=1\) 时，任一“\(1\sqcup(k-1)\)”等变选层机制都需至少 \(k\) 个寄存器态，即至少 \(\log_2 k\) 比特。
\end{theorem}
\begin{proof}
由 \(|A_r|=a\) 可定义 \(\varphi(r):=A_r\in\binom{\Delta}{a}\)。
\(A\) 的不变性等价于
\(
\varphi(\sigma\cdot r)=\sigma\cdot\varphi(r)
\)，故 \(\varphi\) 等变，且 \(A\) 由 \(\varphi\) 唯一恢复，得第 (1) 条。

\(\binom{\Delta}{a}\) 在 \(\mathfrak S_k\) 作用下传递，故非空等变像必为全体 \(\binom{\Delta}{a}\)。
于是
\(
|R|\ge |\mathrm{Im}\,\varphi|=\binom{k}{a}
\)，得第 (2) 条。
第 (3) 条由恒等映射构造直接成立。
最后一句取 \(a=1\) 即得。证毕。
\end{proof}

\begin{theorem}[最小寄存器态下等变 \texorpdfstring{\(a\)}{a}-分裂机制的严格唯一性]\label{thm:conclusion-equivariant-a-splitting-rigidity}
沿定理 \ref{thm:conclusion-equivariant-a-splitting-register-lower-bound} 的记号，定义标准机制
\[
R_{\mathrm{can}}:=\binom{\Delta}{a},
\qquad
A_{\mathrm{can}}
:=
\{(x,B)\in \Delta\times R_{\mathrm{can}}:\ x\in B\}.
\]
若 \((R,A)\) 为一 \(\mathfrak S_k\)-等变 \(a\)-分裂机制，且
\[
|R|=\binom{k}{a},
\]
则作为 \(\mathfrak S_k\)-等变分裂机制有同构
\[
(R,A)\cong (R_{\mathrm{can}},A_{\mathrm{can}}).
\]
此外，
\[
\operatorname{Aut}_{\mathfrak S_k}(R_{\mathrm{can}},A_{\mathrm{can}})=1.
\]
\end{theorem}
\begin{proof}
由定理 \ref{thm:conclusion-equivariant-a-splitting-register-lower-bound}，\((R,A)\) 对应于一个 \(\mathfrak S_k\)-等变映射
\[
\varphi:R\to \binom{\Delta}{a},
\qquad
\varphi(r)=A_r.
\]
同一定理的证明已表明 \(\binom{\Delta}{a}\) 在 \(\mathfrak S_k\) 作用下传递，故 \(\varphi\) 的像必为全体 \(\binom{\Delta}{a}\)。
由于
\[
|R|=\left|\binom{\Delta}{a}\right|=\binom{k}{a},
\]
\(\varphi\) 只能是双射，因而给出 \(\mathfrak S_k\)-集同构。
又由定义，
\[
x\in A_r
\quad\Longleftrightarrow\quad
x\in \varphi(r),
\]
故 \(\varphi\) 同时保持分裂关系，从而得到
\[
(R,A)\cong (R_{\mathrm{can}},A_{\mathrm{can}}).
\]
\par
再设 \(\psi\) 为 \((R_{\mathrm{can}},A_{\mathrm{can}})\) 的一个 \(\mathfrak S_k\)-等变自同构。
由于 \(\psi\) 保持分裂关系，对每个 \(B\in \binom{\Delta}{a}\) 与 \(x\in \Delta\) 都有
\[
x\in B
\quad\Longleftrightarrow\quad
x\in \psi(B).
\]
这迫使 \(\psi(B)=B\) 对所有 \(B\) 成立，故 \(\psi=\mathrm{id}\)。
因此
\[
\operatorname{Aut}_{\mathfrak S_k}(R_{\mathrm{can}},A_{\mathrm{can}})=1.
\]
证毕。
\end{proof}

\begin{theorem}[奇基数非平凡二值跳跃对超对称的结构强迫]\label{thm:conclusion-odd-cardinality-nontrivial-jump-forces-super-symmetry}
设
\[
F:\mathbf{FinSet}^{\cong}\to \mathbf{Tors}(\ZZ_2)
\]
为 monoidal 函子，源端张量为不交并 \(\sqcup\)，目标端张量为 torsor 张量 \(\otimes\)。
若要求其与标准交换同构严格按“无额外符号”的对称约束相容，则对任意奇整数 \(d\ge 3\)，限制
\[
F|_{\mathbf{Set}^{\cong}_d}
\]
必为平凡函子。
等价地：一旦在某个奇基数层出现非平凡二值跳跃（由定理 \ref{thm:bdry-z2-jump-functor-uniqueness} 必同构于 \(\mathrm{Or}\)），交换约束必引入 Koszul 因子，故体系只能是超对称 monoidal，而非无符号对称 monoidal。
\end{theorem}
\begin{proof}
反设存在奇 \(d\ge 3\) 使 \(F|_{\mathbf{Set}^{\cong}_d}\) 非平凡。
由定理 \ref{thm:bdry-z2-jump-functor-uniqueness}，该限制对应符号类。
取 \(|S|=|T|=d\)。
则在块对角子群 \(\mathfrak S_d\times\mathfrak S_d\subset \mathfrak S_{2d}\) 上，\(F\) 的字符限制为两块符号之和；故 \(F\) 在 \(\mathfrak S_{2d}\) 上亦为非平凡字符。
由唯一性只能是 \(\mathrm{sgn}_{2d}\)。
对块交换 \(\beta_{S,T}\)，其奇偶为 \((-1)^{d^2}=-1\)，因此 \(F(\beta_{S,T})\) 必非平凡。

另一方面，若交换约束无额外符号，则当 \(S=T\) 时，\(F(\beta_{S,S})\) 在
\(
F(S)\otimes F(S)
\)
上仅为对换。
由命题 \ref{prop:conclusion-z2torsor-square-parity-collapse} 的规范同构
\(
F(S)\otimes F(S)\cong \ZZ_2
\)，该对换作用为恒等，故 \(F(\beta_{S,S})\) 平凡，矛盾。
故假设不成立。证毕。
\end{proof}
